\documentclass[11pt]{article}
\usepackage[margin=0.82in]{geometry}
\usepackage{amsmath,amssymb,booktabs,longtable,array,graphicx,float}
\usepackage[dvipsnames]{xcolor}
\usepackage{microtype}
\usepackage{enumitem}
\usepackage{caption}
\usepackage{subcaption}
\usepackage{hyperref}
\usepackage{fancyhdr}
\usepackage{tabularx}
\usepackage{makecell}
\usepackage{pdflscape}
\usepackage{lastpage}

\hypersetup{colorlinks=true,linkcolor=MidnightBlue,urlcolor=MidnightBlue,citecolor=MidnightBlue}
\graphicspath{{../../empirical_sid/runs/sid_tier1_primary_20260714_v4/results/final_figures/}}
\setlist{nosep,leftmargin=1.4em}
\captionsetup{font=small,labelfont=bf}
\pagestyle{fancy}
\setlength{\headheight}{14pt}
\setlength{\emergencystretch}{2em}
\fancyhf{}
\lhead{SID empirical validation}
\rhead{14 July 2026}
\cfoot{\thepage\ of \pageref{LastPage}}
\newcommand{\NRMSE}{\mathrm{NRMSE}}
\newcommand{\localtarget}{D_vT(P_h)}
\newcommand{\finitetarget}{\Delta_{T,\delta}(h)}
\newcommand{\SID}{\mathrm{SID}}
\newcommand{\yes}{\textcolor{ForestGreen}{Yes}}
\newcommand{\no}{\textcolor{BrickRed}{No}}
\newcommand{\partialresult}{\textcolor{BurntOrange}{Partial}}
\newcommand{\artifact}[1]{\nolinkurl{#1}}
\IfFileExists{generated_results.tex}{\newcommand{\TierOneMetricRows}{51,000}
\newcommand{\TierTwoMetricRows}{129,480}
\newcommand{\TierOneValidationHash}{\texttt{f537f7917adfe10e...}}
\newcommand{\TierTwoValidationHash}{\texttt{c854efe8cfdf5bb1...}}
}{}

\title{\textbf{Empirical Validation of Score-Derivative SID}\\
\large Local versus finite targets, matched-model controls, dynamic lag recovery, and robustness}
\author{Automated reproducible analysis of the \texttt{new\_sbtg\_neuro} project}
\date{14 July 2026}

\begin{document}
\maketitle

\begin{abstract}
This report executes and audits the staged validation plan in
\emph{SID\_EMPIRICAL\_VALIDATION\_RUNBOOK.md}.  The decisive result is not that one
method wins everywhere.  Correctly specified normalized laws recover every tested
static target and a matched stochastic-gain mixture recovers dynamic distributional
lags, but these models are supplied the relevant mechanism family.  Among reusable
observational estimators, direct moment regression is the most dependable: it
recovers mean, covariance, several third-cumulant constructions, and dynamic
variance/covariance/third-cumulant lags.  Clean score matching sometimes has lower
NRMSE, but its effects are often attenuated by roughly one half.  Multi-noise and
anchored denoising score routes do not provide a general advantage; the anchored
route fails dynamic third cumulants and produces a large observation-only false
effect.  The central bottleneck is the derivative from response score to history
tangent and typed readout, not merely predictive-law fit.  Results support a
targeted predictive-closure program, with score-derived SID retained as a
best-case or reusable-law research direction rather than the current default.
\end{abstract}

\section*{Technical summary}

\noindent\colorbox{gray!10}{\parbox{0.97\linewidth}{
\textbf{Bottom line.}  Use direct targeted regression for a known functional;
use a flexible normalized law when the same fitted law must answer many queries;
use a contrast-aligned classifier for a supported finite intervention when its
representation can express the relevant geometry.  The present score-derived SID
implementations are not competitive with those baselines on calibrated typed
effects, although clean Hyv\"arinen plus Hodge is a useful diagnostic and the
anchored score improves with very large samples.
}}

\begin{itemize}
\item \textbf{Validated artifacts.}  Tier 1 contains 270 static and 30 dynamic
top-level DGP-seed cases, each with a metric table and compressed raw array archive.
All 300 archive hashes match, all primary cells contain 30 common seeds, and no fit
failed.  Tier 2 adds 270 hidden-state/observation-stress cases and 450 static
sample/information cases on new seeds 3001--3030.  A separately frozen post-hoc
fixed-amplitude diagnostic adds 270 cases on paired seeds 4001--4030.
\item \textbf{Matched models are upper bounds, not evidence of generality.}  The
static normalized model knows the signed-tilt basis and fits essentially one
amplitude.  The dynamic ``correct MDN'' knows the gain vector, component noise, and
family, fitting only the mixing law.
\item \textbf{Derivative gap.}  Response-score NRMSE is often 0.04--0.30 while the
corresponding reconstructed history tangent is at or above one, especially for
quartic, tail, and occupancy mechanisms.  Differentiation, centering, and Hodge
integration amplify small response-score errors.
\item \textbf{Dynamic result.}  On observed stochastic gain, matched mixture NRMSE
is about 0.05 and direct moment regression about 0.22--0.26 across variance,
covariance, and third cumulant.  The anchored score reaches variance NRMSE about
0.50 but is slightly under-calibrated, and third-cumulant NRMSE exceeds 1.3 with a
wrong-sign slope.
\item \textbf{Robustness result.}  With a hidden modulator, direct and Gaussian
log-variance recovery improve smoothly with sample size.  The anchored score is
unstable at 8,000 samples and becomes usable only at 32,000, still behind direct
baselines.  Under an observation-only null its false variance effect is comparable
to the true active clean effect.
\item \textbf{Status of the broad claim.}  The evidence is a rigorous diagnostic
validation of these implementations and DGP families, not a validated claim that
SID recovers arbitrary moments in unknown high-dimensional nonlinear systems.
Dense unknown response subspaces, multiple gain factors, and general neural
predictive backends remain unresolved.
\end{itemize}

\section{Question and estimands}

The scientific question is whether a predictive-law model can reveal how a source
history direction changes a chosen future-response functional, particularly when
the conditional mean is invariant but variance, covariance, skewness, tails, or
mode occupancy change.  Two targets must remain separate.

\subsection{Local target}

For conditional law $P_h$ and history direction $v$,
\[
  \ell_v(y;h)=D_v\log p(y\mid h),\qquad
  \localtarget = E_{P_h}\!\left[\mathrm{IF}_T(Y;P_h)\ell_v(Y;h)\right].
\]
The response score $s_y=\nabla_y\log p(y\mid h)$ is not $\ell_v$.  A score route
first estimates $s_y$, differentiates it with respect to history to obtain a mixed
field, reconstructs a centered history tangent (when topology permits), and only
then integrates it against an influence function.

\subsection{Finite target}

For a supported central contrast,
\[
  \finitetarget = \frac{T(P_{h+\delta v})-T(P_{h-\delta v})}{2\delta}.
\]
A balanced endpoint classifier estimates a density-ratio witness for this finite
signed measure.  This is not interchangeable with the local derivative, and the
leaderboards in this report never mix them.

\subsection{Recovery and advantage rules}

The preregistered numerical recovery rule requires an upper 95\% DGP-bootstrap
bound below $\NRMSE=1$, a calibration slope in $[0.8,1.2]$, correct signs, and an
acceptable nuisance false-call rate.  A superiority claim additionally requires
the upper paired error-ratio bound below 0.90 under the same information access.
Because the Tier 1 target-only cases do not emit every nuisance functional, this
report distinguishes \emph{numerical recovery} from the stronger fully validated
recovery label.

\section{Methods map}

\begin{table}[H]
\centering\small
\caption{The four layers that earlier reports sometimes conflated.}
\begin{tabularx}{\linewidth}{p{0.16\linewidth}p{0.25\linewidth}X p{0.18\linewidth}}
\toprule
Layer & Examples & Object learned or transformed & Primary failure mode\\
\midrule
Predictive backend & raw-moment ridge, Gaussian likelihood, matched tilt, matched
gain mixture & $P(Y\mid H)$, selected moments, or structural parameters &
misspecification; target variance\\
Score estimator & clean Hyv\"arinen, multi-noise DSM, Gaussian DSM, anchored
residual energy & $\nabla_y\log p_\sigma(y\mid h)$ & smoothing bias; derivative
instability\\
SID adapter & density AD, ratio critic, energy centering, score secant, Hodge &
$D_v\log p$ or a finite density-ratio witness & centering; support topology;
access mismatch\\
Readout & mean, covariance, cumulants, tail, occupancy, lag tensor & typed
$D_vT$ or $\Delta_{T,\delta}$ & influence-function variance; basis mismatch\\
\bottomrule
\end{tabularx}
\end{table}

\subsection{Methods tested in the new confirmatory suite}

\begin{description}[style=nextline]
\item[Direct raw moments (B2).] Polynomial history basis with heads for raw moments
through order four (and event indicators where required), analytically converted to
central functionals and differentiated.
\item[Matched normalized law (S1/A1).] For M3--M5, the true signed-tilt basis
$\psi(y)$ is supplied and only its amplitude is fitted.  For mean and covariance,
Gaussian moment regressions are used.  This is a specification upper bound.
\item[Ratio critic (A2).] Logistic joint-versus-shuffled classifier with explicit
history--response interactions; the history derivative of its logit is centered
and integrated against a typed influence function.
\item[Clean Hyv\"arinen plus Hodge (S2/A7).] Polynomial response-score model fit by
clean score matching; its history derivative is integrated over response space and
centered.
\item[Multi-noise DSM plus Hodge (S5/A7).] Polynomial denoising score fit at
$\sigma\in\{0.03,0.06,0.12,0.25\}$.
\item[Gaussian DSM control (S6).] Conditional Gaussian score obtained from fitted
low-order moments; this diagnoses how much is available from a Gaussian core.
\item[Anchored residual-energy DSM (S7).] Multi-noise residual score around the
Gaussian core, with development-only ridge and noise ladder frozen before
confirmation.
\item[Finite adapters.] Direct endpoint Monte Carlo, finite differences of moment
regression, normalized-law endpoints, endpoint classifier, signed Riesz dictionary,
and a score-integrated endpoint density.
\end{description}

\section{Reproduction audit}

The frozen legacy core contains 784 completed cases and the finite panel contains
212,832 valid metric rows.  A new nine-case subset crossing G1, G3, and G8 with
Gaussian likelihood, EDM diffusion, and ratio critic was rerun.  All 117 scientific
metric values matched bitwise; 18 runtime-efficiency values varied with execution.
The historical source-tree hash and current source-tree hash differ, so this is
reported as a compatibility reproduction rather than exact source reproduction.

The G8 contrast-aligned developmental experiment remains important: across four
independent generator parameterizations (eight generator/data cells), a motif MLP
has median typed-effect relative error 1.13\%, a typed quartic classifier 1.16\%,
and a full-path/history MLP 3.05\%.  Linear and quadratic negative controls recover
essentially none of the effect.  Thus the G8 failure of generic generative readouts
is not non-identifiability; it is representation/readout misalignment.

\IfFileExists{../../empirical_sid/runs/sid_tier1_primary_20260714_v4/results/final_figures/g8_typed_classifier.pdf}{
\begin{figure}[H]\centering
\includegraphics[width=0.82\linewidth]{g8_typed_classifier.pdf}
\caption{G8 typed finite contrast: aligned classifiers recover the quartic path
effect while lower-order negative controls do not.  This is developmental evidence
with four independent systems, not a 30-system confirmation.}
\end{figure}}{}

\section{Oracle gates}

Stage 0 executed 53 algebra/software checks covering normalization, positivity,
moment orthogonality, local and finite identities, Gaussian corruption, dynamic lag
derivatives, classifier conversion, and Hodge reconstruction.  All passed.  The
topology ladder is decisive: Hodge reconstruction is numerically exact while bridge
mass is positive, but a pure disconnected mixture-weight change is globally
unrecoverable from within-component response gradients.

\begin{figure}[H]\centering
\includegraphics[width=0.70\linewidth]{topology_hodge.pdf}
\caption{Oracle Hodge reconstruction across the registered support-topology ladder.
The disconnected endpoint is an identifiability failure, not an optimizer failure.}
\end{figure}

The oracle gates establish that the requested typed quantities are mathematically
identifiable in the connected static systems and that the numerical readout formulas
are correct.  They do not establish that a learned score supplies an accurate mixed
field.

\section{Score-estimator results}

\begin{figure}[H]\centering
\includegraphics[width=\linewidth]{derivative_gap_heatmap.pdf}
\caption{Response-score and reconstructed history-tangent NRMSE use the same fitted
score models.  Accurate response scores frequently fail after history
differentiation and Hodge reconstruction.  A star denotes an upper 95\% bound below
one.}
\end{figure}

The clean Hyv\"arinen estimator is the strongest generic score route in the static
suite.  On the M3 variants its typed NRMSE is roughly 0.54--0.57, but calibration
slopes are only about 0.47--0.49.  The apparent NRMSE win over direct moments is
therefore shrinkage, not calibrated recovery.  Multi-noise DSM is consistently
worse, and the Gaussian and anchored controls remain near zero-effect behavior on
M3 while exploding on fourth-order targets.

The derivative gap supports the following failure chain:
\[
\text{small score error}
\ \not\Rightarrow\ 
\text{small mixed-field error}
\ \not\Rightarrow\ 
\text{small centered-tangent error}
\ \not\Rightarrow\ 
\text{calibrated typed effect}.
\]
Stabilization reduces numerical variance but imposes a strong low-order/zero-effect
bias.  More denoising noise does not repair this structural mismatch.

\section{Local SID leaderboard}

\begin{figure}[H]\centering
\includegraphics[width=\linewidth]{static_local_heatmap.pdf}
\caption{Mean local NRMSE over 30 top-level DGP seeds.  The matched normalized law
is separated visually because it is supplied the DGP tilt basis.  Stars denote an
upper 95\% bound below one; they do not by themselves certify calibration or
nuisance specificity.}
\end{figure}

\IfFileExists{tables/local_key_results.tex}{\begin{longtable}{p{0.15\linewidth}p{0.23\linewidth}p{0.23\linewidth}p{0.23\linewidth}c}
\caption{Selected local results. Values are mean [95\% DGP-bootstrap interval]. Partial means NRMSE recovery without registered calibration.}\\
\toprule Target & Method & NRMSE & Calibration slope & Calibrated?\\ \midrule
\endfirsthead \toprule Target & Method & NRMSE & Calibration slope & Calibrated?\\ \midrule \endhead
M1 mean & Direct moments & 0.515 [0.421, 0.622] & 1.010 [0.955, 1.063] & \yes \\
M1 mean & Ratio critic & 0.447 [0.391, 0.511] & 1.008 [0.931, 1.082] & \yes \\
M1 mean & Clean Hyv\"arinen + Hodge & 0.324 [0.290, 0.361] & 0.836 [0.794, 0.875] & \yes \\
M2 covariance & Direct moments & 0.453 [0.383, 0.526] & 1.021 [0.968, 1.073] & \yes \\
M2 covariance & Ratio critic & 0.506 [0.446, 0.570] & 1.050 [0.954, 1.151] & \yes \\
M2 covariance & Gaussian score control & 0.453 [0.385, 0.527] & 1.021 [0.968, 1.071] & \yes \\
M3 bounded skew & Matched normalized law & 0.105 [0.080, 0.133] & 1.008 [0.961, 1.055] & \yes \\
M3 bounded skew & Direct moments & 0.658 [0.562, 0.763] & 0.990 [0.916, 1.063] & \yes \\
M3 bounded skew & Ratio critic & 0.718 [0.620, 0.825] & 1.003 [0.888, 1.126] & \yes \\
M3 bounded skew & Clean Hyv\"arinen + Hodge & 0.547 [0.515, 0.579] & 0.489 [0.450, 0.530] & \partialresult \\
M4 quartic & Matched normalized law & 0.088 [0.069, 0.110] & 0.993 [0.956, 1.031] & \yes \\
M4 quartic & Direct moments & 4.443 [3.900, 5.041] & 0.684 [0.194, 1.182] & \no \\
M4 quartic & Clean Hyv\"arinen + Hodge & 3.970 [3.697, 4.249] & 0.099 [-0.223, 0.419] & \no \\
M4 tail & Matched normalized law & 0.193 [0.099, 0.296] & 0.807 [0.713, 0.895] & \yes \\
M4 tail & Direct moments & 8.278 [6.962, 9.735] & 1.025 [0.071, 1.947] & \no \\
M5 occupancy & Matched normalized law & 0.084 [0.051, 0.119] & 0.916 [0.882, 0.950] & \yes \\
M5 occupancy & Direct moments & 2.569 [2.071, 3.112] & 0.897 [0.636, 1.172] & \no \\
\bottomrule\end{longtable}
}{}

The general observational conclusion is heterogeneous.  Direct moments recover
M1 mean, M2 covariance, and all four M3 third-cumulant constructions numerically,
with near-unit slopes.  The ratio critic is competitive on M1/M2 and two M3
variants but needs conditional evaluation draws in this implementation.  No
general method recovers M4 quartic at medium information.  M4 tail and M5 occupancy
are additionally information-bound because positivity caps achieved
$\Lambda$ at 0.128 and 0.949, respectively.  Their matched-model success proves the
targets exist in the data but does not remove the information limitation for a
generic learner.

\section{Finite leaderboard}

\begin{figure}[H]\centering
\includegraphics[width=\linewidth]{static_finite_heatmap.pdf}
\caption{Finite-target NRMSE at $\delta=0.12$.  Local and finite values are never
pooled.  Endpoint labels or queries are stronger information access than
observational regression.}
\end{figure}

Finite differences of fitted moment regression behave similarly to the local
direct method.  Naive independent endpoint Monte Carlo is extremely noisy because
two noisy functional estimates are divided by $2\delta$; common random numbers or
analytic model moments are necessary.  The generic endpoint classifier and signed
Riesz dictionary do not recover most static channels at this sample size.  This
does not contradict G8: the successful G8 classifier uses a contrast-aligned
quartic/motif representation and far more evaluation draws.

\section{Higher-order findings}

The experiments distinguish four ideas that should not be called simply
``higher-order recovery.''

\begin{enumerate}
\item \textbf{Literal M3 skew.}  Raw cubic moment heads recover several signed-tilt
constructions; clean score matching finds the direction but attenuates it.
\item \textbf{M4 fourth central moment.}  Every general route fails at medium
information, including direct fourth raw moments.  The matched normalized tilt
recovers it, demonstrating estimation variance/basis mismatch rather than
non-identifiability.
\item \textbf{Tail and occupancy.}  These are event probabilities, not merely
polynomial order.  Positivity prevents medium information in the registered family,
so their negative result is a detection-boundary result as well as an adapter
failure.
\item \textbf{G8 path feature.}  The effect lives in a structured whole-path motif.
An aligned classifier succeeds while generic law samplers and lower-order
classifiers miss it.  Geometry and representation dominate nominal moment order.
\end{enumerate}

\section{Dynamic lag results}

\begin{figure}[H]\centering
\includegraphics[width=0.93\linewidth]{dynamic_nrmse.pdf}
\caption{D1/D2 local lag-profile NRMSE with DGP-bootstrap intervals.  The matched
gain mixture is a specification upper bound; direct moments are the strongest
general observational baseline.}
\end{figure}

\begin{figure}[H]\centering
\includegraphics[width=\linewidth]{dynamic_lag_profiles.pdf}
\caption{Recovered source-to-target lag profiles, averaged across 30 DGP seeds.
The score route follows variance imperfectly and reverses/loses the third-cumulant
profile, whereas direct moments and the matched mixture preserve shape and sign.}
\end{figure}

On the D1 fixed-variance mean anchor, polynomial mean regression and Gaussian
likelihood tie at NRMSE about 0.049, while the ratio critic is about 0.164.  This is
the expected no-advantage result for a correctly specified mean problem.

On D2, the matched mixture achieves NRMSE 0.054--0.057.  Direct moments achieve
0.221 for covariance, 0.258 for variance, and 0.262 for third cumulant, all with
near-unit slope and sign accuracy above 0.99.  The ratio critic is usable but worse.
The anchored score has variance NRMSE 0.496 but slope 0.797, just outside the
registered interval, and third-cumulant NRMSE 1.329 with a negative slope.  Original
SBTG field norm is reported only as an unsigned localization diagnostic and is not
reinterpreted as a tangent.

\section{Baseline comparison and information access}

\IfFileExists{tables/information_access.tex}{\begin{table}[H]\centering\small
\caption{Audited information access.}
\begin{tabularx}{\linewidth}{p{0.22\linewidth}p{0.27\linewidth}X}\toprule
Method class & Effective access & Consequence\\\midrule
Direct moments & observational training & Target-specific raw moment/event heads; analytic derivatives \\
Ratio critic & training plus conditional evaluation queries & Fresh conditional draws center and integrate the witness \\
Generic score readouts (Tier 1) & training plus oracle-density quadrature & Favorable best case; not an observational-only comparison \\
Generic score readouts (Tier 2 static) & training plus oracle-density quadrature & Same favorable quadrature audit as Tier 1 \\
Matched normalized law & known DGP tilt basis & Fits the amplitude of supplied $\psi(y)$ \\
Matched dynamic mixture & known gain/noise family & Fits the mixing law with supplied structural form \\
Endpoint adapters & endpoint labels or conditional queries & Stronger access than observational local regression \\
\bottomrule\end{tabularx}\end{table}
}{}

Several nominal method labels in the raw runner understate their access.  The audit
therefore preserves the recorded label and adds an audited label:

\begin{itemize}
\item static and Tier 1 dynamic score readouts use true-DGP density quadrature for
centering/integration, making their results favorable best-case score controls;
\item ratio readouts use fresh conditional evaluation samples;
\item the static matched law knows $\psi(y)$, and the dynamic matched mixture knows
gain/noise structure;
\item direct moment regression uses observational training and analytic model
derivatives.
\end{itemize}

Consequently no score or matched-model result demonstrates observational advantage
over direct moments.  The matched-model results instead quantify the benefit of
correct structural knowledge.

\section{Failure analysis}

\begin{table}[H]
\centering\small
\caption{Where each route fails.}
\begin{tabularx}{\linewidth}{p{0.22\linewidth}p{0.22\linewidth}X}
\toprule
Route & First failing stage & Evidence and interpretation\\
\midrule
Direct moments & readout variance/basis & M4 quartic NRMSE $>4$; tail and occupancy
fail at low achieved information; works well for M1--M3 and D2\\
Ratio critic & tangent/readout variance & competitive on simple laws, degrades with
context and high-order targets; requires conditional evaluation queries here\\
Clean Hyv\"arinen & calibration after Hodge & M3 NRMSE looks favorable but slope is
about one half; fourth-order effects fail\\
Multi-noise DSM & mixed-field derivative & response score is tolerable, history
tangent and typed effects are poor\\
Anchored score & stabilization bias and derivative variance & shrinks M3 toward
zero, fails D2 third cumulant, unstable under hidden context, false effect under
observation-only null\\
Endpoint classifier/Riesz & witness variance/representation & generic static
adapters fail; aligned G8 classifier succeeds\\
Matched normalized laws & external validity & excellent in-family, but supplied the
mechanism family and therefore not a generic solution\\
\bottomrule
\end{tabularx}
\end{table}

\section{Robustness and support}

\begin{figure}[H]\centering
\includegraphics[width=\linewidth]{tier2_dynamic.pdf}
\caption{Tier 2 hidden-modulator and observation-filter stresses on new seeds
3001--3030.  Direct/log-variance routes scale predictably; the anchored score is
sample hungry and sensitive to observation filtering.}
\end{figure}

\begin{figure}[H]\centering
\includegraphics[width=\linewidth]{tier2_static_slices.pdf}
\caption{Sample-size and achieved-information slices for bounded M3, M4 quartic,
and M5 occupancy.  The top row is a preregistered local-alternative sequence:
target $\Lambda=3$ is held fixed by recalibrating amplitude as $N$ changes.  The
middle row is the post-hoc ordinary learning curve with amplitude fixed at its
$N=8{,}000$ value.  The bottom row varies target information at $N=8{,}000$.
Positivity bounds make some target values unattainable; actual achieved
$\Lambda$ is used and stored.  The design is not a Cartesian product.}
\end{figure}

\IfFileExists{tables/tier2_key_results.tex}{\begin{longtable}{p{0.17\linewidth}p{0.18\linewidth}p{0.08\linewidth}p{0.20\linewidth}p{0.22\linewidth}}
\caption{Selected Tier 2 results: registered rows use seeds 3001--3030; post-hoc fixed-amplitude rows use seeds 4001--4030.}\\
\toprule Cell & Method & Metric & Mean [95\% CI] & Audited access\\\midrule
\endfirsthead\toprule Cell & Method & Metric & Mean [95\% CI] & Audited access\\\midrule\endhead
D3 context 2, N=8k & Gaussian log-variance & nrmse & 0.183 [0.152, 0.214] & observational \\
D3 context 2, N=8k & Direct moments & nrmse & 0.236 [0.198, 0.272] & observational \\
D3 context 2, N=8k & Anchored score + Hodge & nrmse & 0.383 [0.308, 0.457] & observational training plus conditional evaluation queries \\
D3 context 12, N=32k & Gaussian log-variance & nrmse & 0.178 [0.157, 0.200] & observational \\
D3 context 12, N=32k & Direct moments & nrmse & 0.212 [0.195, 0.229] & observational \\
D3 context 12, N=32k & Anchored score + Hodge & nrmse & 0.397 [0.373, 0.424] & observational training plus conditional evaluation queries \\
D5 primary filter & Direct moments & nrmse & 0.331 [0.302, 0.359] & observational \\
D5 primary filter & Anchored score + Hodge & nrmse & 0.531 [0.480, 0.580] & observational training plus conditional evaluation queries \\
D5 observation-only null & Anchored score + Hodge & null\_rms & 0.079 [0.064, 0.102] & observational training plus conditional evaluation queries \\
M3 bounded, N=2k & Direct moments & NRMSE & 0.972 [0.846, 1.108] & observational \\
M3 bounded, N=32k & Direct moments & NRMSE & 0.628 [0.527, 0.732] & observational \\
M4 quartic, N=32k & Direct moments & NRMSE & 3.709 [3.279, 4.138] & observational \\
M4 quartic, N=32k & Matched normalized law & NRMSE & 0.079 [0.058, 0.101] & known dgp tilt basis plus fitted amplitude \\
M5 occupancy, N=32k & Direct moments & NRMSE & 1.843 [1.612, 2.086] & observational \\
M5 occupancy, N=32k & Matched normalized law & NRMSE & 0.109 [0.089, 0.130] & known dgp tilt basis plus fitted amplitude \\
Post-hoc M3 fixed amp., N=2k & Direct moments & NRMSE & 1.158 [0.982, 1.335] & observational \\
Post-hoc M3 fixed amp., N=32k & Direct moments & NRMSE & 0.284 [0.247, 0.323] & observational \\
Post-hoc M4 fixed amp., N=2k & Direct moments & NRMSE & 6.850 [5.594, 8.135] & observational \\
Post-hoc M4 fixed amp., N=32k & Direct moments & NRMSE & 1.943 [1.665, 2.232] & observational \\
Post-hoc M4 score, N=2k & Clean Hyv\"arinen + Hodge & NRMSE & 4.679 [4.078, 5.413] & observational training plus oracle density quadrature \\
Post-hoc M4 score, N=32k & Clean Hyv\"arinen + Hodge & NRMSE & 4.066 [3.888, 4.238] & observational training plus oracle density quadrature \\
Post-hoc M5 fixed amp., N=2k & Direct moments & NRMSE & 4.312 [3.756, 4.881] & observational \\
Post-hoc M5 fixed amp., N=32k & Direct moments & NRMSE & 1.474 [1.227, 1.762] & observational \\
\bottomrule\end{longtable}
}{}

The hidden-modulator experiment tests predictive closure: short context is easier
because fewer lag coordinates are requested, sufficient context recovers the full
12-lag source history, and overcomplete context adds eight true null lags.  At
$N=32{,}000$ the anchored score can recover variance numerically, but it remains
less efficient than direct/log-variance estimators.  This leaves open whether a
more appropriate neural score architecture can close the gap; it does not support
the current anchored polynomial implementation as the preferred route.

The static sample-size row is not an ordinary fixed-effect learning curve.  It
holds total local information near $\Lambda=3$, so the perturbation amplitude
shrinks as $N$ grows.  It asks whether each estimator remains efficient under
contiguous local alternatives.  The separate information row holds
$N=8{,}000$ and varies target $\Lambda$; positivity bounds prevent event families
from reaching every requested regime.

The separately frozen post-hoc fixed-amplitude curve resolves the ambiguity.
For M3 bounded skew, direct-moment NRMSE improves from 1.16 at $N=2{,}000$ to
0.56 at $8{,}000$ and 0.28 at $32{,}000$, with slope near one throughout; the
ratio critic improves from 1.12 to 0.69 to 0.45.  Clean score/Hodge stays near
0.56 NRMSE at all three sizes and its slope stays near 0.46--0.48.  More data
therefore repairs sampling error for the direct/ratio routes but reveals an
attenuation plateau in the score route.  Because this diagnostic was motivated
after the fixed-$\Lambda$ results were inspected, it is explanatory rather than
confirmatory.

M4 quartic is different.  At fixed amplitude, direct-moment NRMSE falls from
6.85 to 3.77 to 1.94 as $N$ increases, so sampling error matters, but it still
does not beat the zero estimator at $32{,}000$ (achieved $\Lambda\approx12$).
Clean score/Hodge changes only from 4.68 to 4.38 to 4.07.  The matched normalized
law improves from 0.22 to 0.12 to 0.06.  Thus the quartic target is learnable and
additional data help the direct estimator, but the registered generic routes are
far too inefficient and the score route shows another error plateau.

M5 occupancy improves without crossing the recovery threshold.  At fixed
amplitude, direct-event NRMSE falls from 4.31 to 2.88 to 1.47, the ratio critic
from 4.37 to 2.27 to 1.27, and clean score/Hodge from 2.53 to 1.76 to 1.55.
The matched law falls from 0.16 to 0.08 to 0.05.  At $32{,}000$, achieved
$\Lambda\approx3.8$ is no longer intrinsically tiny, so the remaining failure is
not explained by the original $N=8{,}000$ positivity boundary alone.

The observation filter stress is more damaging.  Filtering spreads an active
variance effect across adjacent source lags, which direct moments still recover.
In the matched observation-only null, the score route emits a false variance
profile comparable in RMS to the active clean effect.  Any real neural-data claim
therefore requires an explicit observation-process null and deconvolution or
measurement model.

\section{Claim registry}

\IfFileExists{tables/claim_registry.tex}{\begin{table}[H]\centering\small
\caption{Preregistered claim registry.}
\begin{tabularx}{\linewidth}{p{0.05\linewidth}X p{0.10\linewidth}p{0.37\linewidth}}\toprule
ID & Claim & Result & Evidence\\\midrule
H1 & Response-score accuracy need not imply history-tangent accuracy & \yes & Paired operator panels show large error amplification \\
H2 & Direct history tangent can recover typed effects & \partialresult & Supported for M1--M3 and D1--D3; not M4/tail/occupancy \\
H3 & Stabilized score route has a general advantage & \no & No equal-access paired advantage; calibration and null failures \\
H4 & Hodge recovery depends on connected support & \yes & Oracle bridge ladder passes until disconnected endpoint \\
H5 & Finite contrasts can be easier than local derivatives & \partialresult & True for G8 aligned representation; not generic endpoint adapters \\
H6 & Higher-order success depends on geometry/representation & \yes & M3, M4, events, and path motif separate sharply \\
H7 & Distributional lag recovery is possible & \yes & Matched mixture and direct moments recover D2 profiles \\
\bottomrule\end{tabularx}\end{table}
}{}

The evidence supports the derivative-gap hypothesis, calibrated direct recovery
for several typed targets, topology dependence of Hodge reconstruction, and
geometry dependence of higher-order recovery.  It rejects a general advantage for
the tested stabilized score route and rejects treating original SBTG magnitude as
a signed typed tangent.  It leaves unresolved generic dense high-dimensional SID,
unknown affected subspaces, multi-source/target low-rank factors, and foundation
model transfer.

\section{Decision}

\IfFileExists{tables/final_decision.tex}{\begin{landscape}\scriptsize\begin{longtable}{p{0.08\linewidth}p{0.12\linewidth}p{0.12\linewidth}p{0.12\linewidth}p{0.12\linewidth}p{0.08\linewidth}p{0.08\linewidth}p{0.15\linewidth}}
\caption{Final target-by-target decision table. NRMSE values are mean [95\% DGP-bootstrap interval].}\\
\toprule Target & Best direct & Best normalized law & Best score-derived & Best finite & Recovered? & Advantage? & Main limitation\\\midrule
\endfirsthead\toprule Target & Best direct & Best normalized law & Best score-derived & Best finite & Recovered? & Advantage? & Main limitation\\\midrule\endhead
Mean & Direct moments 0.52 [0.42, 0.62] & Matched normalized law 0.52 [0.42, 0.62] & Clean Hyv\"arinen + Hodge 0.32 [0.29, 0.36] & Direct moments (finite) 0.50 [0.41, 0.61] & Numerically yes & No score advantage & Low-order home field; full nuisance FPR not emitted \\
Variance lag & Direct moments 0.26 [0.23, 0.29] & Matched gain mixture 0.05 [0.05, 0.06] & Anchored score + Hodge 0.50 [0.44, 0.55] & Direct moments (finite) 0.26 [0.23, 0.29] & Yes: direct/matched & Matched only, with stronger structure & Score slope 0.797; hidden state and filtering add error \\
Covariance lag & Direct moments 0.22 [0.20, 0.25] & Matched gain mixture 0.05 [0.05, 0.06] & No score readout & Direct moments (finite) 0.22 [0.20, 0.25] & Yes: direct/matched & Matched only, with stronger structure & No generic score covariance estimate in the implemented D2 route \\
Third-cumulant lag & Direct moments 0.26 [0.22, 0.31] & Matched gain mixture 0.06 [0.05, 0.06] & Anchored score + Hodge 1.33 [1.27, 1.38] & Direct moments (finite) 0.26 [0.22, 0.31] & Yes: direct/matched & Matched only, with stronger structure & Score slope is negative; static clean score is attenuated \\
Standardized skewness & Direct raw moments available & Matched-law moments available & Not separately confirmed & Endpoint moments available & Unresolved as a primary claim & No & Run emitted third cumulant, not a standalone primary skewness board \\
Tail probability & Direct moments 8.28 [6.96, 9.73] & Matched normalized law 0.19 [0.10, 0.30] & Clean Hyv\"arinen + Hodge 6.26 [5.63, 6.92] & Direct moments (finite) 8.01 [6.66, 9.39] & Matched law only & No fair general advantage & Positivity capped achieved $\Lambda$ at 0.128 \\
Mode occupancy & Direct moments 2.57 [2.07, 3.11] & Matched normalized law 0.08 [0.05, 0.12] & Clean Hyv\"arinen + Hodge 1.85 [1.70, 2.00] & Direct moments (finite) 2.48 [2.01, 2.98] & Matched law only & No fair general advantage & Boundary-limited; achieved $\Lambda$ at target 3 is 0.949 \\
Fourth central / quartic path & Direct moments 4.44 [3.90, 5.04] & Matched normalized law 0.09 [0.07, 0.11] & Clean Hyv\"arinen + Hodge 3.97 [3.70, 4.25] & G8 aligned motif MLP: 1.13\% median relative error (development) & Matched/access-aligned only & No general advantage & G8 is finite and developmental; it is not the M4 local target \\
Dynamic lag tensor & Direct moments: 0.22--0.26 across D2 distributional channels & Matched gain mixture: 0.054--0.057 & 0.50 variance; 1.33 third cumulant & Direct finite: 0.22--0.26 & Yes: direct/matched & No score advantage & Dense unknown subspaces and multiple gain factors remain untested \\
\bottomrule\end{longtable}\end{landscape}
\begin{table}[H]\centering\small
\caption{Method-selection decision from the implemented evidence.}
\begin{tabularx}{\linewidth}{p{0.28\linewidth}p{0.28\linewidth}X}\toprule
Use case & Preferred route & Reason\\\midrule
Known mean/covariance/cumulant/event target & Targeted direct head & Best general calibration and efficiency in implemented cells \\
Many queries from one reusable conditional law & Flexible normalized MDN/flow & Promising next baseline; current matched laws are specification upper bounds \\
Supported finite intervention with known geometry & Contrast-aligned classifier & G8 development evidence; generic classifiers are insufficient \\
Operator diagnosis or theory study & Clean score + mixed-field/Hodge gates & Useful decomposition, not current production estimator \\
Primary general SID estimator & Do not select current DSM/anchored score & No calibrated/equal-access advantage and observation-null failure \\
\bottomrule\end{tabularx}\end{table}
}{}

\paragraph{Recommended scientific direction.}
Build the next method around \emph{targeted predictive closure}: a shared history
encoder with calibrated heads for means, covariance factors, selected cumulants,
tail/occupancy logits, and possibly characteristic-function or kernel features.
Benchmark it against a flexible normalized MDN/flow when amortized multi-query use
justifies learning the full law.  Retain ratio critics for supported contrasts and
clean score/Hodge as an operator diagnostic.  Do not make the diffusion/DSM score
route the primary estimator until it passes mixed-field, calibration, and
observation-null gates.

\paragraph{What would change this decision.}
A convincing score-derived SID result would need: (i) a generic score backend with
accurate history mixed fields on at least two independent high-order DGP families;
(ii) calibrated typed effects with no oracle density quadrature; (iii) paired
advantage over direct and normalized-law baselines under equal access and compute;
(iv) robustness to hidden state and observation filters; and (v) dense unknown
response geometry rather than a supplied affected coordinate.

\section{Limitations and next experiments}

\begin{enumerate}
\item \textbf{Dense D4 benchmark.}  Rotate several independent gain factors into an
unknown response subspace, use multiple source histories and timescales, and score
the full source-by-target-by-lag tensor before low-rank factorization.
\item \textbf{Fair flexible laws.}  Compare shared-encoder direct heads with a
capacity-matched MDN and conditional flow; the current ``correct MDN'' is an upper
bound, not a flexible baseline.
\item \textbf{Score architecture.}  Replace the compact polynomial score with a
time-series network, but gate it first on clean score, mixed-field, and tangent
truth.  Explore explicit Jacobian regularization, integrability penalties, and
joint history-response score formulations.
\item \textbf{Readout without oracle quadrature.}  Estimate centering and influence
functions from held-out observational samples or backend samples, propagate that
uncertainty, and use common evaluation budgets.
\item \textbf{Finite contrast efficiency.}  Use coupled/common-random-number
sampling, calibrated density-ratio classifiers, conditional MMD, and analytic
feature dictionaries where scientifically justified.
\item \textbf{Observation model.}  Add calcium/voltage filtering, missingness,
heteroskedastic noise, and a matched observation-only null before real-neural-data
interpretation.
\item \textbf{Mechanism confusion.}  Emit every readout on every DGP so the
preregistered nuisance false-call requirement is directly estimable rather than
inferred from target-only cases.
\item \textbf{Unfinished runbook controls.}  Complete standalone standardized
skewness, the full $\delta/\sigma$ ladders, teacher-contamination and stabilization
ablations, heavy-tail/manifold stresses, and multiple stochastic neural
initializations.  The compact estimators here are deterministic given a DGP seed
and therefore do not substitute for those neural-backend tests.
\end{enumerate}

\section*{Further questions}

\begin{itemize}
\item Is the reusable object really the full response score, or a learned basis of
influence functions/conditional characteristic functions that supports many typed
queries more directly?
\item Can the history tangent be learned directly with normalization/centering
constraints, avoiding differentiation of a response score?
\item Which neuromodulatory questions require local derivatives, and which are more
naturally supported finite interventions?
\item How should measurement dynamics and biological slow gain be jointly
identified when both create lag spreading?
\end{itemize}

\appendix
\section{Negative results and invalidated executions}

Negative results were retained rather than removed from the boards.  The principal
scientific failures are: all general M4 quartic routes at medium information;
all general tail and occupancy routes in their boundary-limited registered
families; generic static endpoint classifiers and signed Riesz witnesses; the
multi-noise DSM history tangent; anchored-score D2 third cumulants; and the
anchored-score observation-only null.  No failed confirmatory fit was silently
dropped: both Tier 1 and Tier 2 validated with zero failed top-level cases.

\begin{table}[H]\centering\small
\caption{Developmental or invalidated executions excluded from confirmation.}
\begin{tabularx}{\linewidth}{p{0.18\linewidth}p{0.24\linewidth}X}
\toprule Run/version & Status & Reason and disposition\\\midrule
Tier 1 v1 & Development only & Positive-noise score fields were centered against
the wrong law; amendment 1 corrected the declared noised object before
confirmation.\\
Anchored tuning grid & Development only & One global ridge/noise setting was
selected on eight development seeds; its poor development performance was retained
and the setting frozen.\\
Tier 1 v3 & Invalidated after 9/270 cells & Raw arrays were not persisted.  The run
was stopped before results were inspected; v4 used identical scientific settings
and new immutable CSV/NPZ pairs.\\
Initial legacy audit launch & Rejected before data & The legacy tier enumeration
rejected \texttt{audit\_reproduction}; it was changed to the valid
\texttt{development} label without changing scientific settings.\\
Legacy rerun & Compatibility only & All 117 scientific values were bitwise equal,
but source-tree hashes differ and runtime fields are environment-dependent.\\
\bottomrule
\end{tabularx}
\end{table}

The full amendment trail is preserved in
\artifact{empirical_sid/PREREGISTRATION_AMENDMENT_1.md} through
\artifact{PREREGISTRATION_AMENDMENT_3.md} and
\artifact{LEGACY_AUDIT_AMENDMENT_20260714.md}.

\section{DGP and implementation details}

The static primary suite uses response dimension eight, 8,000 training examples,
96 evaluation histories, $\delta=0.12$, and 30 top-level seeds.  M3--M5 effects lie
in a registered scalar projection with seven independent nuisance coordinates.
Signed tilts are positive normalized laws
\[
  p_h(y)=\phi(y)\{1+a\tanh(h)\psi(y)\},
\]
with $\psi$ orthogonalized against lower-order nuisance polynomials.  This design
isolates functional order but does not test discovery of an unknown dense response
subspace.

The static direct baseline uses ridge regression on the scalar history basis
$1,h,\ldots,h^5$ with simultaneous raw-response heads through order four and an
additional event-indicator head for tail/occupancy channels.  Central-functional
derivatives are then computed analytically, including all mean-correction terms.
The ratio critic is a joint-versus-shuffled logistic regression using cubic history
features, response powers through five, smooth/local event features, and every
history--response interaction; its logit derivative is centered with fresh
conditional draws.  These are deliberately compact estimators, not neural
foundation models.

The score models use a tensor-product polynomial basis through cubic history and
quintic scaled response, with a noise-coordinate block.  Clean Hyv\"arinen solves
the empirical score-matching normal equations.  Multi-noise DSM uses corrupted
samples at $\sigma\in\{0.03,0.06,0.12,0.25\}$; the anchored version subtracts a
Gaussian score constructed from fitted mean/variance heads and regularizes the
residual coefficients.  The mixed field is a central history secant, and in the
one-dimensional affected coordinate Hodge reconstruction is numerical integration
followed by mean-zero centering.  For the Tier 1 and static Tier 2 score boards,
that centering/influence integration uses true-DGP density quadrature, intentionally
giving the score route a favorable upper-control readout.

D2 uses a Bernoulli stochastic-gain state centered to preserve conditional mean:
\[
  Y=a(S-\pi(h))+\epsilon,
  \qquad \pi(h)=\operatorname{logit}^{-1}(\alpha+k^\top h).
\]
Its variance, covariance, and third-cumulant lag derivatives are analytic.  D3
integrates a logistic-normal hidden modulator with Gauss--Hermite quadrature.  D5
filters independent adjacent latent responses and adds observation noise; its
variance derivative is the sum of current and shifted filtered contributions.

\section{Statistical analysis}

All primary intervals resample the top-level DGP seed 2,000 times.  Nested query
draws are averaged within seed.  No confirmatory result selected a hyperparameter;
anchored-score ridge and noise schedule were frozen on eight development seeds.
Paired comparisons use common DGP seeds.  Test outcomes are never pooled with
development runs or the aborted pre-raw-array execution.

\IfFileExists{tables/compute_summary.tex}{\begin{table}[H]\centering\small
\caption{Recorded sequential compute. Wall time is summed once per DGP-family/seed cell; peak memory is the maximum runner-recorded resident usage.}
\begin{tabular}{lrrrr}\toprule
Component & Cells & Total CPU-wall hours & Median sec/cell & Peak MiB\\\midrule
Tier 1 static & 270 & 0.43 & 7.15 & 306 \\
Tier 1 dynamic & 60 & 0.01 & 0.50 & n/a \\
Tier 2 dynamic robustness & 270 & 0.07 & 0.30 & n/a \\
Tier 2 static slices & 450 & 0.86 & 6.81 & 580 \\
Post-hoc fixed-amplitude & 270 & 0.51 & 6.76 & 627 \\
\bottomrule\end{tabular}\end{table}
}{}

All runs were executed sequentially on CPU with one BLAS/OpenMP thread.  The
Tier 2 static process ran at reduced operating-system priority; completion checks
also enforced paired CSV/NPZ counts, SHA-256 agreement, finite metrics, and a
minimum 2\,GiB free-disk stop condition.  The machine-observed process resident
set remained below 0.7\,GiB.

\section{Artifact and reproducibility ledger}

\IfFileExists{tables/artifact_ledger.tex}{\scriptsize\begin{longtable}{p{0.17\linewidth}p{0.56\linewidth}p{0.16\linewidth}}
\caption{Primary artifact ledger. The displayed hashes are prefixes; complete hashes are in package manifests.}\\
\toprule Role & Path & SHA-256 prefix\\\midrule
\endfirsthead\toprule Role & Path & SHA-256 prefix\\\midrule\endhead
Tier 1 summary & \artifact{/Users/vik/Developer/new_sbtg_neuro/empirical_sid/runs/sid_tier1_primary_20260714_v4/results/summary.csv} & \texttt{30dba5e0dcaf...} \\
Tier 1 validation & \artifact{/Users/vik/Developer/new_sbtg_neuro/empirical_sid/runs/sid_tier1_primary_20260714_v4/results/validation.json} & \texttt{f537f7917adf...} \\
Tier 1 package manifest & \artifact{/Users/vik/Developer/new_sbtg_neuro/empirical_sid/runs/sid_tier1_primary_20260714_v4/results/package_manifest.json} & \texttt{4bae88517c22...} \\
Tier 2 summary & \artifact{/Users/vik/Developer/new_sbtg_neuro/empirical_sid/runs/sid_tier1_primary_20260714_v4/results/tier2/summary.csv} & \texttt{d0fca3e239a2...} \\
Tier 2 validation & \artifact{/Users/vik/Developer/new_sbtg_neuro/empirical_sid/runs/sid_tier1_primary_20260714_v4/results/tier2/validation.json} & \texttt{c854efe8cfdf...} \\
Tier 2 package manifest & \artifact{/Users/vik/Developer/new_sbtg_neuro/empirical_sid/runs/sid_tier1_primary_20260714_v4/results/tier2/package_manifest.json} & \texttt{386827fcc276...} \\
Legacy audit & \artifact{/Users/vik/Developer/new_sbtg_neuro/empirical_sid/runs/sid_tier1_primary_20260714_v4/reports/legacy_audit/legacy_reproduction_audit.json} & \texttt{1c6224b5e886...} \\
Post-hoc diagnostic plan & \artifact{/Users/vik/Developer/new_sbtg_neuro/empirical_sid/TIER2_POSTHOC_FIXED_AMPLITUDE_PLAN_20260714.md} & \texttt{9f6e0e709b0a...} \\
Machine-readable claims registry & \artifact{/Users/vik/Developer/new_sbtg_neuro/empirical_sid/runs/sid_tier1_primary_20260714_v4/results/claims_registry.csv} & \texttt{a4ebe4490b3b...} \\
\bottomrule\end{longtable}
}{}

Primary machine-readable outputs are \artifact{results/seed_level.parquet},
\artifact{results/summary.csv}, \artifact{results/tables/paired_advantage.csv},
the Tier 2 package, raw per-seed NPZ archives, manifests, frozen configuration, and
source snapshots.  SHA-256 inventories and validation JSON files accompany both
tiers.

\section{Runbook completion checklist}

\IfFileExists{tables/completion_checklist.tex}{\begin{table}[H]\centering\small
\caption{Runbook completion checklist.}
\begin{tabularx}{\linewidth}{X p{0.10\linewidth}p{0.42\linewidth}}\toprule
Item & Complete? & Audit note\\\midrule
Frozen preregistration/config/source & \yes & Hashes recorded before confirmatory execution \\
Stage 0 algebra/software gates & \yes & 53 checks passed \\
Tier 1 static and dynamic cases & \yes & 270 static + 30 dynamic; 30 common seeds \\
Legacy compatibility reproduction & \yes & 117 scientific metrics bitwise equal; source hashes differ \\
Tier 2 hidden-state/observation stresses & \yes & 270 dynamic + 450 static registered cases \\
Post-hoc fixed-amplitude diagnostic & \yes & 270 exploratory cases on new seeds; cannot upgrade confirmation \\
DGP-seed bootstrap and access audit & \yes & 2,000 replicates; effective access retained \\
Full learned mechanism-confusion matrix & \no & Target-only cells cannot certify nuisance FPR \\
Dense unknown response geometry (D4) & \no & Affected coordinate is supplied in current static suite \\
Flexible MDN/flow and neural score backends & \no & Still proposed work \\
Full stabilization/teacher-contamination ablations & \no & Frozen compact score shortlist only \\
Standalone skewness and full delta/noise ladders & \no & Primary registered cells only \\
\bottomrule\end{tabularx}\end{table}
}{}

The program is complete for the implemented Tier 1 and registered Tier 2 cells.
The broad runbook completion claim remains withheld because full learned
mechanism-confusion matrices, dense D4 geometry, flexible MDN/flow robustness, and
general neural score backends were not completed.  The scientifically correct
label is therefore \textbf{diagnostic empirical validation}, not universal
moment-recovery validation.

\end{document}
